\section{Find Square Root of a Number using Binary Search}
\label{sec:bs-sqrt}

\noindent We now begin the second major group of problems in this chapter:
\textbf{binary search on the answer}. In every problem so far, the
candidate space was a set of array indices. From here on, there is often no
array at all -- the candidate space is a range of numbers that could be the
\emph{answer itself}, and we binary search directly over that range,
checking feasibility of each candidate with a helper function.

\problemstatement{Given a non-negative integer $n$, find $\lfloor
\sqrt{n} \rfloor$, the floor of its square root, without using a built-in
square root function -- and ideally without testing every integer from $1$
to $n$ one at a time.}

\begin{patternbox}
\emph{``Find the floor of the square root''} doesn't mention an array at
all, which is exactly the point: the phrase to watch for is \emph{``find
the largest/smallest value such that some computable condition holds.''}
Here, the condition is $\textit{candidate}^2 \le n$. This condition is
monotonic in \textit{candidate} -- true for small candidates, false for
large ones, with exactly one crossover -- which is all binary search ever
needed in the first place. The array, in earlier sections, was just one
particular way of presenting a monotonic candidate space; an integer range
works just as well.
\end{patternbox}

\subsection*{Understanding the problem carefully}

Define $P(\textit{candidate}) = (\textit{candidate}^2 \le n)$ for
candidates ranging over $0, 1, 2, \dots, n$. As \textit{candidate}
increases, $\textit{candidate}^2$ increases too (squaring is monotonic for
non-negative numbers), so $P$ is \texttt{true} for every candidate up to
$\lfloor \sqrt{n} \rfloor$ and \texttt{false} for every candidate beyond
it. The floor of the square root is exactly the \emph{largest} candidate
for which $P$ holds.

\begin{ideabox}[Candidate Space and Predicate]
The candidate space is the integer range $[0, n]$ (or, more tightly,
$[1, n]$ for $n \ge 1$). The predicate is
$P(\textit{candidate}) = (\textit{candidate}^2 \le n)$, and we want the
largest candidate where $P$ holds -- the same ``find the rightmost true
position'' template used for floor in Section~\ref{sec:bs-floor-ceil},
just applied to a numeric range instead of array values.
\end{ideabox}

\subsection*{Why binary searching the answer directly is valid}

There is nothing conceptually new here beyond recognizing that the
candidate space need not be array indices. The monotonicity of $P$ is
guaranteed by basic algebra ($x \le y \Rightarrow x^2 \le y^2$ for
non-negative $x, y$), which plays exactly the role that array sortedness
played in earlier sections: it guarantees a single, well-defined crossover
point that binary search can locate in logarithmically many evaluations of
$P$, each evaluation costing only $O(1)$ (a single multiplication).

\subsection*{Algorithm}

\begin{enumerate}
  \item Initialize $\textit{lo} \gets 0$, $\textit{hi} \gets n$,
        $\textit{ans} \gets 0$.
  \item While $\textit{lo} \le \textit{hi}$:
  \begin{enumerate}
    \item $\textit{mid} \gets \textit{lo}+(\textit{hi}-\textit{lo})/2$.
    \item If $\textit{mid} \times \textit{mid} \le n$: record
          $\textit{ans} \gets \textit{mid}$; search right:
          $\textit{lo} \gets \textit{mid}+1$.
    \item Else: search left: $\textit{hi} \gets \textit{mid}-1$.
  \end{enumerate}
  \item Return \textit{ans}.
\end{enumerate}

\begin{algorithm}[H]
\caption{Floor of Square Root}
\begin{algorithmic}[1]
\Procedure{FloorSqrt}{$n$}
  \State $\textit{lo} \gets 0$, $\textit{hi} \gets n$, $\textit{ans} \gets 0$
  \While{$\textit{lo} \le \textit{hi}$}
    \State $\textit{mid} \gets \textit{lo}+(\textit{hi}-\textit{lo})/2$
    \If{$\textit{mid} \times \textit{mid} \le n$}
      \State $\textit{ans} \gets \textit{mid}$
      \State $\textit{lo} \gets \textit{mid} + 1$
    \Else
      \State $\textit{hi} \gets \textit{mid} - 1$
    \EndIf
  \EndWhile
  \State \Return \textit{ans}
\EndProcedure
\end{algorithmic}
\end{algorithm}

\subsection*{Dry run}

\dryrun{Take $n = 28$.}

\begin{itemize}
  \item $\textit{lo}=0,\textit{hi}=28$: $\textit{mid}=14$, $14^2=196 \le
        28$? No. $\textit{hi}=13$.
  \item $\textit{lo}=0,\textit{hi}=13$: $\textit{mid}=6$, $6^2=36 \le 28$?
        No. $\textit{hi}=5$.
  \item $\textit{lo}=0,\textit{hi}=5$: $\textit{mid}=2$, $2^2=4 \le 28$?
        Yes. $\textit{ans}=2$, $\textit{lo}=3$.
  \item $\textit{lo}=3,\textit{hi}=5$: $\textit{mid}=4$, $4^2=16 \le 28$?
        Yes. $\textit{ans}=4$, $\textit{lo}=5$.
  \item $\textit{lo}=5,\textit{hi}=5$: $\textit{mid}=5$, $5^2=25 \le 28$?
        Yes. $\textit{ans}=5$, $\textit{lo}=6$.
  \item $\textit{lo}=6 > \textit{hi}=5$: loop ends.
\end{itemize}

The answer is $5$, and indeed $5^2 = 25 \le 28 < 36 = 6^2$.

\subsection*{C++ implementation}

\begin{lstlisting}
#include <bits/stdc++.h>
using namespace std;

int floorSqrt(int n) {
    int lo = 0, hi = n, ans = 0;

    while (lo <= hi) {
        long long mid = lo + (hi - lo) / 2;
        if (mid * mid <= n) {
            ans = mid;
            lo = mid + 1;
        } else {
            hi = mid - 1;
        }
    }
    return ans;
}

int main() {
    cout << "Floor sqrt of 28: " << floorSqrt(28) << endl;
    return 0;
}
\end{lstlisting}

\begin{complexitybox}
\textbf{Time Complexity.} The candidate range $[0, n]$ halves each
iteration, giving
\[
  O(\log n).
\]
\textbf{Space Complexity.} $O(1)$ auxiliary space.
\end{complexitybox}

\begin{takeawaybox}
Binary search on the answer requires exactly two ingredients: a numeric
candidate range, and a feasibility predicate that is monotonic across that
range. Once both exist, the same \texttt{while(lo <= hi)} skeleton from
array-based binary search applies verbatim -- only the meaning of
\textit{lo}, \textit{hi}, and the predicate check change.
\end{takeawaybox}
